\documentclass[fontsize=15pt]{jlreq}

%プリアンブル
\usepackage{amsmath}
\usepackage{mathtools}
\pagestyle{empty}
\usepackage[top=15mm, bottom=20mm, left=20mm, right=20mm]{geometry}

\newcommand{\m}[1]{\raisebox{0.1em}{\textcircled{\scriptsize #1}}}



%本文
\begin{document}
\begin{center}
{\LARGE {振り返り}} \\
-生物\m{2}- %単元ごとに変えるか消す
\end{center}
\vspace{1em}
\begin{flushright}
名前：\underline{\hspace{5cm}}
\end{flushright}

% 問題部分
\begin{enumerate}
    \item 親がもつ形質の内，どちらか一方しか現れない互いに対をなす形質を何というか．\\
    (\hspace{4cm})
    \vspace{2em} 

    \item 生物のからだをつくる細胞の染色体の内，同じ形質を示す，2本ずつ対をなしている染色体のことをなんというか？\\
    (\hspace{4cm})
    \vspace{2em} 

    \item 対立形質のそれぞれについての純系同士を交配させたとき，\m{1}子に現れる形質と，\m{2}子に現れない形質をそれぞれなんというか．また，エンドウの種についての例を挙げよ．\\
    \m{1}(\hspace{3cm})(例\hspace{4cm})\\
    \m{2}(\hspace{3cm})(例\hspace{4cm})
    
    \vspace{2em} 

    \item 親から子，子から孫に発現する対立形質の割合について，メンデルの実験から分かることを簡潔に述べよ．
    \begin{itemize}
        \item 親から子：
        \item 子から孫：
    \end{itemize}
    \vspace{2em} 

    \item 染色体に含まれる物質であり，遺伝子の本体を何と言うか日本語で答えよ．(アルファベット×)\\
    (\hspace{4cm})
    \vspace{2em} 

\end{enumerate}

\end{document}
